

\section*{\SGAInlineRefTarget{sga3:ii:section:4}4. Tangent Space to a Group --- Lie Algebras}
\addcontentsline{toc}{section}{4. Tangent Space to a Group --- Lie Algebras}

\subsection*{\SGAInlineRefTarget{sga3:ii:subsection:4-1}4.1.}
Let \(G\) be a group functor over \(S\). We have seen that
\(T_{G/S}(\mathcal M)\) and
\(\underline{\operatorname{Lie}}(G/S,\mathcal M)\) are then endowed with
group structures over \(S\). There are group morphisms (3.1)

\SGAThreeNativeDiagram[0.52\textwidth]
  {native_diagrams/exp2/exp2_p022_split_tangent_sequence.tex}
  {The split tangent sequence of the group functor \(G\).}
  {fig:exp2-split-tangent-sequence}
  {LNM 151, pp.~59--60}

By definition, \(i\) identifies
\(\underline{\operatorname{Lie}}(G/S,\mathcal M)\) with the kernel of
\(p\), and \(s\) is a section of \(p\). It follows from \SGARefLink{sga3:i:proposition:2-3-7}{I~2.3.7} that this
sequence identifies \(T_{G/S}(\mathcal M)\) with the semidirect product
of \(G\) by \(\underline{\operatorname{Lie}}(G/S,\mathcal M)\).
\SGARefTarget{sga3:ii:definition:4-1-a}%
The corresponding action of \(G\) on
\(\underline{\operatorname{Lie}}(G/S,\mathcal M)\) is denoted
\(\operatorname{Ad}\) and is called the \emph{adjoint representation} of
\(G\); thus, by definition, for \(x\in G(S')\) and
\(X\in\underline{\operatorname{Lie}}(G/S,\mathcal M)(S')\),
\[
  \operatorname{Ad}(x)X
  =i^{-1}\bigl(s(x)i(X)s(x)^{-1}\bigr).
\]
If \(G\) is commutative, then \(T_{G/S}(\mathcal M)\) is commutative as
well and \(\operatorname{Ad}(x)X=X\).

\SGARefTarget{sga3:ii:definition:4-1-b}%
Let \(G,H\) be group functors over \(S\), and let \(f:G\to H\) be a group
morphism. Functoriality gives a morphism of exact sequences compatible
with the canonical sections.

\SGAThreeNativeDiagram[0.66\textwidth]
  {native_diagrams/exp2/exp2_p022_derived_exact_sequences.tex}
  {Functoriality of the split tangent sequence.}
  {fig:exp2-derived-exact-sequences}
  {LNM 151, pp.~59--60}

The morphism \(L(f)\), also denoted
\(\underline{\operatorname{Lie}}(f)\), is called the \emph{derived
morphism} of \(f\).\footnote{Translator's note: the Gille--Polo re-edition
(2011), Expos\'e~II, adds that if \(G/S\) and \(H/S\) satisfy (E), then
\(L(f)\) is \(\mathcal O_S\)-linear for the module structures of 3.6; this
is 3.7~bis applied to the unit sections, since \(f(e)=e\).}

\SGAInlineRefTarget{sga3:ii:proposition:4-1-1}\subsection*{\SGARefLink{sga3:ii:proposition:4-1-1}{Proposition 4.1.1}.}
Let \(g\in G(S)\). Then
\[
  \operatorname{Ad}(g):
  \underline{\operatorname{Lie}}(G/S,\mathcal M)
  \longrightarrow
  \underline{\operatorname{Lie}}(G/S,\mathcal M)
\]
is the derived morphism of
\(\operatorname{Int}(g):G\to G\).

Indeed,
\(\operatorname{Ad}(g)X=i^{-1}(\operatorname{Int}(g)i(X))\), which is
precisely \(L(\operatorname{Int}(g))(X)\) by the definition of the
derived morphism.

\SGAInlineRefTarget{sga3:ii:corollary:4-1-1-1}\subsection*{Corollary.}
Suppose that \(G/S\) satisfies (E). Then we know
(\SGARefLink{sga3:ii:proposition:3-9}{3.9})\footnote{Translator's note: the 1970 text refers here to
``3.10 cor.~1''; the statement used is 3.9 (or Corollary~3 after 3.10), to
which the Gille--Polo re-edition (2011), Expos\'e~II, refers. The
re-edition also invokes the functoriality of the operators of
\(\mathcal O_S\): \(\operatorname{Ad}(g)=L(\operatorname{Int}(g))\) is
\(\mathcal O_S\)-linear by 3.7~bis.} that the group structure of
\(\underline{\operatorname{Lie}}(G/S,\mathcal M)\) defined above is none
other than the structure induced by its \(\mathcal O_S\)-module structure
(defined by means of (E)). It follows from the preceding proposition that
the \(\operatorname{Ad}(g)\) respect the \(\mathcal O_S\)-module structure
of \(\underline{\operatorname{Lie}}(G/S,\mathcal M)\):
\[
  \operatorname{Ad}:G\longrightarrow
  \underline{\operatorname{Aut}}_{\mathcal O_S\text{-mod.}}
    \bigl(\underline{\operatorname{Lie}}(G/S,\mathcal M)\bigr);
\]
in other words, \(\operatorname{Ad}\) is a \emph{linear representation} of
\(G\) in the \(\mathcal O_S\)-module
\(\underline{\operatorname{Lie}}(G/S,\mathcal M)\).

\SGAInlineRefTarget{sga3:ii:remark:4-1-1-2}\subsection*{Remarks.}
\textit{1)} The functor \(G/S\) satisfies \(E\) if and only if, for every
pair \((\mathcal M,\mathcal N)\) of finite free
\(\mathcal O_S\)-modules, the following diagram, obtained by applying
\(\underline{\operatorname{Lie}}(G/S,-)\) to diagram \((*)\) of \SGARefLink{sga3:ii:definition:2-1}{2.1}, is
cartesian.

\SGAThreeNativeDiagram[0.45\textwidth]
  {native_diagrams/exp2/exp2_p023_condition_e_cartesian.tex}
  {The cartesian criterion for condition \(E\).}
  {fig:exp2-condition-e-cartesian}
  {LNM 151, p.~60}

\textit{2)} If \(G/S\) satisfies \(E\), the derived morphism of the group
law \(\pi:G\times_SG\to G\) is none other than the addition law in
\(\underline{\operatorname{Lie}}(G/S,\mathcal M)\).

\textit{Nota.} \(\pi\) is not a group morphism, but \(\pi(e,e)=e\), which
suffices to define \(L(\pi)\) (cf. \SGARefLink{sga3:ii:proposition:3-7}{3.7}).\footnote{Translator's note: the
Gille--Polo re-edition (2011), Expos\'e~II, adds that \(L(\pi)\) maps
\(L^{(e,e)}_{G\times_SG/S}(\mathcal M)=\underline{\operatorname{Lie}}(G/S,\mathcal M)
\times_S\underline{\operatorname{Lie}}(G/S,\mathcal M)\) (3.8) to
\(\underline{\operatorname{Lie}}(G/S,\mathcal M)\), and that likewise the
derived morphism of \(g\mapsto g^n\) is multiplication by \(n\) (see the
Addition after Corollary~3 of 3.10).}

Let us now study the set \(\Gamma(T_{G/S}(\mathcal M)/G)\). Note first
that there is an isomorphism
\[
  \operatorname{Hom}\bigl(G,\underline{\operatorname{Lie}}(G/S,\mathcal M)\bigr)
  \xrightarrow{\sim}\Gamma\bigl(T_{G/S}(\mathcal M)/G\bigr)
\]
defined as follows: to
\(f:G\to\underline{\operatorname{Lie}}(G/S,\mathcal M)\) one associates
the section \(s_f:G\to T_{G/S}(\mathcal M)\) such that
\[
  s_f(g)=i(f(g))s(g),\qquad g\in G(S'),\ S'\to S.
\]
Let \(h\) be an automorphism of the functor \(G\) over \(S\), not
necessarily respecting the group structure. To every section \(t\) of
\(T_{G/S}(\mathcal M)\) one can associate \(h(t)\), defined by transport
of structure: it is, for example, the unique section of
\(T_{G/S}(\mathcal M)\) that makes the following diagram commute.

\SGAThreeNativeDiagram[0.43\textwidth]
  {native_diagrams/exp2/exp2_p024_transport_square.tex}
  {Transport of a tangent section by an automorphism \(h\).}
  {fig:exp2-transport-square}
  {LNM 151, pp.~61--62}

In particular, take for \(h\) the right translation by an element \(x\) of
\(G(S)\):
\[
  h(g)=t_x(g)=g\cdot x,\qquad g\in G(S'),\ S'\to S.
\]
One has at once
\[
  t_x(s_f)=s_{t_x(f)},
\]
where \(t_x(f)\) denotes the morphism of \(G\) into
\(\underline{\operatorname{Lie}}(G/S,\mathcal M)\) defined by
\[
  t_x(f)(g)=f(g\cdot x^{-1}),\qquad g\in G(S'),\ S'\to S.
\]
It follows that if \(G\) is made to operate by right translations on
\(\Gamma(T_{G/S}(\mathcal M)/G)\) and on
\(\operatorname{Hom}(G,\underline{\operatorname{Lie}}(G/S,\mathcal M))\),
the isomorphism defined above respects the operations of \(G\). In
particular, the sections of \(T_{G/S}(\mathcal M)\) invariant under right
translation correspond under this isomorphism to the \emph{constant}
morphisms from \(G\) to \(\underline{\operatorname{Lie}}(G/S,\mathcal M)\)
(that is, those factoring through the projection \(G\to S\)), or again to
the elements of \(\underline{\operatorname{Lie}}(G/S,\mathcal M)(S)\). In
other words:\footnote{Translator's note: the Gille--Polo re-edition (2011),
Expos\'e~II, carries this out functorially in \(S'\): for \(x\in G(S')\)
it lets \(x\) act on
\(\underline{\operatorname{Hom}}_{G/S}(G,T_{G/S}(\mathcal M))\) and on
\(\underline{\operatorname{Hom}}_S(G,\underline{\operatorname{Lie}}(G/S,\mathcal M))\)
by \((\sigma\cdot x)(g)=\sigma(gx^{-1})\,s(x)\) and
\((f\cdot x)(g)=f(gx^{-1})\); the map \(f\mapsto s_f\) is equivariant,
since \(s\) is a group morphism, and the invariant sections are called
sections invariant under right translation. Likewise, with \(x\) acting on
\(\underline{\operatorname{End}}_{I_S(\mathcal M)/S}(G_{I_S(\mathcal M)})\)
by \((u\cdot x)(g)=u(gx^{-1})\,x\), it makes precise that an endomorphism
\(u\) ``commutes with right translations'' in 4.1.3 when \(u\cdot x=u\)
after every base change \(S''\to S\) and for every \(x\in G(S'')\); the map
of the Corollary of 3.12 is equivariant.}

\SGAInlineRefTarget{sga3:ii:proposition:4-1-2}\subsection*{\SGARefLink{sga3:ii:proposition:4-1-2}{Proposition 4.1.2}.}
The map
\[
 \underline{\operatorname{Lie}}(G/S,\mathcal M)(S)
 \longrightarrow \Gamma(T_{G/S}(\mathcal M)/G)
\]
that associates with \(X\in\underline{\operatorname{Lie}}(G/S,\mathcal M)(S)\)
the section \(x\mapsto X\cdot x\) is an isomorphism of
\(\underline{\operatorname{Lie}}(G/S,\mathcal M)(S)\) onto the part of
\(\Gamma(T_{G/S}(\mathcal M)/G)\) formed by the sections invariant under
right translation.
\footnote{Statements 4.1.2, 4.1.3, and 4.1.4 can be obtained more simply
by observing that the automorphisms of \(G\) invariant under right
translations are the left translations.}

Taking \SGARefLink{sga3:ii:corollary:3-12-1}{3.12, Corollary}, into account, one obtains:

\SGAInlineRefTarget{sga3:ii:proposition:4-1-3}\subsection*{\SGARefLink{sga3:ii:proposition:4-1-3}{Proposition 4.1.3}.}
There is an isomorphism, functorial in \(G\), between
\(\underline{\operatorname{Lie}}(G/S,\mathcal M)(S)\) and the set of
\(I_S(\mathcal M)\)-endomorphisms of \(G_{I_S(\mathcal M)}\) that
induce the identity on \(G\) and commute with the right translations of
\(G\).

Taking \SGARefLink{sga3:ii:proposition:3-13}{3.13} into account gives the following theorem.

\SGAInlineRefTarget{sga3:ii:theorem:4-1-4}\subsection*{\SGARefLink{sga3:ii:theorem:4-1-4}{Theorem 4.1.4}.}
Let \(G\) be a functor in groups over \(S\) satisfying (E) relative to
\(S\). The group \(\underline{\operatorname{Lie}}(G/S,\mathcal M)(S)\)
identifies, functorially in \(G\), with the group of
\(I_S(\mathcal M)\)-automorphisms of \(G_{I_S(\mathcal M)}\) that
induce the identity on \(G\) and commute with right translations.

For \(\mathcal M=\mathcal O_S\), this recovers one of the classical
definitions of the Lie algebra of a group.

\subsection*{\SGAInlineRefTarget{sga3:ii:subsection:4-2}4.2.}
Before going further, let us establish some new corollaries of
\SGARefLink{sga3:ii:proposition:3-11}{3.11}. We have seen in \emph{loc.\ cit.}\ that there is a canonical
isomorphism
\[
 L^u_{\underline{\operatorname{Hom}}_S(X,Y)/S}(\mathcal M)
 \xrightarrow{\sim}
 \underline{\operatorname{Hom}}_{Y/S}(X,T_{Y/S}(\mathcal M));
\]
suppose that \(Y\) is a group; then
\(T_{Y/S}(\mathcal M)=\underline{\operatorname{Lie}}(Y/S,\mathcal M)\cdot Y
=\underline{\operatorname{Lie}}(Y/S,\mathcal M)_Y\), whence an isomorphism
\[
 L^u_{\underline{\operatorname{Hom}}_S(X,Y)/S}(\mathcal M)
 \xrightarrow{\sim}
 \underline{\operatorname{Hom}}_S
   \bigl(X,\underline{\operatorname{Lie}}(Y/S,\mathcal M)\bigr).
\]

\subsection*{Proposition 4.2.}
Let \(u:X\to Y\) be a morphism of \(S\)-groups. There is a functorial
isomorphism
\[
 L^u_{\underline{\operatorname{Hom}}_{S\text{-gr.}}(X,Y)/S}
   (\mathcal M)
 \xrightarrow{\sim}
 \underline Z^1_S
   \bigl(X,\underline{\operatorname{Lie}}(Y/S,\mathcal M)\bigr).
\]

[\(\underline Z^1_S\) denotes the ``functor of crossed homomorphisms'',
defined in a manner analogous to the functor
\(\underline{\operatorname{Hom}}\). We let \(X\) operate on
\(\underline{\operatorname{Lie}}(Y/S,\mathcal M)\) through the morphism
\[
 X\xrightarrow{\ u\ }Y\xrightarrow{\ \operatorname{ad}\ }
 \underline{\operatorname{Aut}}_{S\text{-gr.}}
   \bigl(\underline{\operatorname{Lie}}(Y/S,\mathcal M)\bigr)\,\bigr].
\]

This follows from the isomorphism
\[
 L^u_{\underline{\operatorname{Hom}}_{S\text{-gr.}}(X,Y)/S}(\mathcal M)
 \xrightarrow{\sim}
 \underline{\operatorname{Hom}}_{(Y/S)\text{-gr.}}
 \bigl(X,\underline{\operatorname{Lie}}(Y/S,\mathcal M)\cdot Y\bigr)
\]
and from the well-known identification of the sections of a semidirect
product with the cocycles of the quotient in the kernel (see for example
III~1.2.2).

\par\medskip\noindent
\emph{Addition} (after the Gille--Polo re-edition (2011), Expos\'e~II; our
formulation). The re-edition makes this proof explicit; in our words:
(a) For \(Z\)-groups \(X\) and \(Y\), let
\(\underline{\operatorname{Hom}}_{(Z/S)\text{-gr.}}(X,Y)\) be the subfunctor
of \(\underline{\operatorname{Hom}}_{Z/S}(X,Y)\) whose value on \(S'\) is
\(\operatorname{Hom}_{Z_{S'}\text{-gr.}}(X_{S'},Y_{S'})\). Under the
isomorphism of 3.11, a point of
\(T_{\underline{\operatorname{Hom}}_{(Z/S)\text{-gr.}}(X,Y)/S}(\mathcal M)\)
is, after base change, a \(Z\)-morphism
\(\varphi\colon X\times_SI_S(\mathcal M)\to Y\) multiplicative in the first
variable, that is, a morphism of \(Z\)-groups \(X\to T_{Y/Z}(\mathcal M)\),
the law of \(T_{Y/Z}(\mathcal M)\) being obtained from that of \(Y\) by
applying \(\underline{\operatorname{Hom}}_Z(I_Z(\mathcal M),-)\). Hence
\(T_{\underline{\operatorname{Hom}}_{(Z/S)\text{-gr.}}(X,Y)/S}(\mathcal M)
\simeq\underline{\operatorname{Hom}}_{(Z/S)\text{-gr.}}(X,T_{Y/Z}(\mathcal M))\),
functorially in \(\mathcal M\); for \(Z=S\) and the fibre over \(u\), this
gives the isomorphism used above, \(T_{Y/S}(\mathcal M)\) being
\(\underline{\operatorname{Lie}}(Y/S,\mathcal M)\cdot Y\). The same holds for
\(\mathcal O_S\)-modules and \(\mathcal O_S\)-linear maps, with the module
structures deduced from that of \(Y\). (b) For an \(S\)-group \(L\) on
which \(X\) acts by group automorphisms, \(\underline Z^1_S(X,L)\) is the
subfunctor of \(\underline{\operatorname{Hom}}_S(X,L)\) formed by the
\(\varphi\) such that \(\varphi(x_1x_2)=\varphi(x_1)\,(x_1\cdot\varphi(x_2))\)
on points. A morphism \(\Phi\colon X\to L\cdot Y\) lying over \(u\) is
written uniquely \(\Phi(x)=\varphi(x)\,u(x)\) with \(\varphi\colon X\to L\),
and \(\Phi\) is a group morphism if and only if
\(\varphi(x_1x_2)=\varphi(x_1)\,u(x_1)\varphi(x_2)u(x_1)^{-1}\), the cocycle
condition for the action through \(\operatorname{ad}\circ u\). (c) One has
\(\underline Z^1_S(X,L\times_SL')\simeq\underline Z^1_S(X,L)\times_S
\underline Z^1_S(X,L')\); if \(L\) is an \(\mathcal O_S\)-module on which
\(X\) acts linearly, \(\underline Z^1_S(X,L)\) is the kernel of the
differential \(\underline{\operatorname{Hom}}_S(X,L)\to
\underline{\operatorname{Hom}}_S(X\times_SX,L)\) of \SGARefLink{sga3:i:subsection:5-1}{I~5.1}, hence an
\(\mathcal O_S\)-module. (d) If \(Y/S\) satisfies (E), so does
\(\underline{\operatorname{Hom}}_{S\text{-gr.}}(X,Y)/S\), by (a), since
\(\underline{\operatorname{Hom}}_{S\text{-gr.}}(X,-)\) commutes with fibre
products; both sides of the proposition then carry \(\mathcal O_S\)-module
structures deduced from functoriality in \(\mathcal M\), which on the right
is the structure of (c) (\(X\) acts linearly by the Corollary of 4.1.1), and
the isomorphism, being functorial in \(\mathcal M\), is
\(\mathcal O_S\)-linear.
\par\medskip

Applying 3.10, Corollary~1, one deduces:

\SGAInlineRefTarget{sga3:ii:corollary:4-2-1}\subsection*{Corollary 1.}
If \(Y/S\) satisfies (E), there are functorial isomorphisms
\[
 L^u_{\underline{\operatorname{Isom}}_{S\text{-gr.}}(X,Y)/S}
   (\mathcal M)
 \xrightarrow{\sim}
 \underline Z^1_S
   \bigl(X,\underline{\operatorname{Lie}}(Y/S,\mathcal M)\bigr).
\]

\SGAInlineRefTarget{sga3:ii:corollary:4-2-2}\subsection*{Corollary 2.}
If \(X/S\) satisfies (E), there are functorial isomorphisms
\[
 \underline{\operatorname{Lie}}
   \bigl(\underline{\operatorname{Aut}}_{S\text{-gr.}}(X)/S,
         \mathcal M\bigr)
 \xrightarrow{\sim}
 \underline Z^1_S
   \bigl(X,\underline{\operatorname{Lie}}(X/S,\mathcal M)\bigr).
\]

\SGAInlineRefTarget{sga3:ii:corollary:4-2-3}\subsection*{Corollary 3.}
If \(Y\) is commutative, there is a functorial
isomorphism\footnote{Translator's note: the Gille--Polo re-edition (2011),
Expos\'e~II, gives the reason: for commutative \(Y\) the adjoint
representation is trivial (4.1), so the crossed homomorphisms for the
action through \(u\) are the group morphisms.}
\[
 L^u_{\underline{\operatorname{Hom}}_{S\text{-gr.}}(X,Y)/S}
   (\mathcal M)
 \xrightarrow{\sim}
 \underline{\operatorname{Hom}}_{S\text{-gr.}}
   \bigl(X,\underline{\operatorname{Lie}}(Y/S,\mathcal M)\bigr).
\]

\subsection*{\SGAInlineRefTarget{sga3:ii:subsection:4-3}4.3.}
Let us now consider the case in which \(X\) and \(Y\) are
\(\mathcal O_S\)-modules. If \(Y\) is an \(\mathcal O_S\)-module, the
functor \(\underline{\operatorname{Lie}}(Y/S,\mathcal M)\) is equipped with
an \(\mathcal O_S\)-module structure deduced from that of \(Y\). Equipped
with this structure, it will provisionally be denoted
\(\underline{\operatorname{Lie}}'(Y/S,\mathcal M)\).\footnote{Translator's note: the prime is used to distinguish the
\(\mathcal O_S\)-module structure induced by that of \(Y\) from the one
that will be obtained from condition \(E\).}

When \(Y/S\) satisfies (E), we shall always write
\(\underline{\operatorname{Lie}}(Y/S,\mathcal M)\) for the functor
\(\underline{\operatorname{Lie}}(Y/S,\mathcal M)\) equipped with the
\(\mathcal O_S\)-module structure defined for every functor satisfying
(E). We know that the abelian-group structures of
\(\underline{\operatorname{Lie}}(Y/S,\mathcal M)\) and
\(\underline{\operatorname{Lie}}'(Y/S,\mathcal M)\) coincide, but a priori
the same is not true of the module structures (see a counterexample in
\SGARefLink{sga3:ii:section:6}{section~6}).

As an immediate consequence of the preceding corollary, one has:

\subsection*{Proposition 4.3.}
There is a functorial isomorphism
\[
 L^u_{\underline{\operatorname{Hom}}_{\mathcal O_S\text{-mod.}}(X,Y)/S}
   (\mathcal M)
 \xrightarrow{\sim}
 \underline{\operatorname{Hom}}_{\mathcal O_S\text{-mod.}}
   \bigl(X,\underline{\operatorname{Lie}}'(Y/S,\mathcal M)\bigr).\footnotemark
\]
\footnotetext{Translator's note: the Gille--Polo re-edition (2011),
Expos\'e~II, adds that if \(Y/S\) satisfies (E), so does
\(\underline{\operatorname{Hom}}_{\mathcal O_S\text{-mod.}}(X,Y)/S\), and the
isomorphism is then \(\mathcal O_S\)-linear for the structures deduced from
functoriality in \(\mathcal M\); on the right this is
\((a\varphi)(x)=a\cdot\varphi(x)\), with the operation of (E) on
\(\underline{\operatorname{Lie}}(Y/S,\mathcal M)\), which differs from
\((a\varphi)(x)=\varphi(ax)\) unless the two structures of 4.3 agree. It
also notes that, through \(\varphi\mapsto u\oplus\varphi\), the isomorphism
is a restriction of \(T_{\underline{\operatorname{Hom}}_{\mathcal O_S\text{-mod.}}(X,Y)/S}(\mathcal M)\simeq
\underline{\operatorname{Hom}}_{\mathcal O_S\text{-mod.}}(X,Y\oplus\underline{\operatorname{Lie}}'(Y/S,\mathcal M))\),
and that, when \(Y/S\) satisfies (E) and \(u\) is an isomorphism, the
tangent spaces at \(u\) of \(\underline{\operatorname{Isom}}_{\mathcal O_S\text{-mod.}}(X,Y)\)
and of \(\underline{\operatorname{Hom}}_{\mathcal O_S\text{-mod.}}(X,Y)\)
coincide, by Corollary~1 of 3.13.}

\SGAInlineRefTarget{sga3:ii:corollary:4-3-1}\subsection*{Corollary.}
If \(X/S\) satisfies (E), there are functorial isomorphisms
\[
 \underline{\operatorname{Lie}}
  \bigl(\underline{\operatorname{Aut}}_{\mathcal O_S\text{-mod.}}(X)/S,
        \mathcal M\bigr)
 \xrightarrow{\sim}
 \underline{\operatorname{Hom}}_{\mathcal O_S\text{-mod.}}
  \bigl(X,\underline{\operatorname{Lie}}'(X/S,\mathcal M)\bigr).\footnotemark
\]
\footnotetext{Translator's note: the Gille--Polo re-edition (2011),
Expos\'e~II, adds that this isomorphism respects the structures deduced
from functoriality in \(\mathcal M\), and that
\(\underline{\operatorname{Aut}}_{\mathcal O_S\text{-mod.}}(X)/S\) satisfies
(E): as \(X/S\) satisfies (E),
\(\underline{\operatorname{Lie}}'(X/S,\mathcal M\oplus\mathcal N)\simeq
\underline{\operatorname{Lie}}'(X/S,\mathcal M)\times_S
\underline{\operatorname{Lie}}'(X/S,\mathcal N)\), so both sides carry
direct sums to products, which suffices by Remark~1) of 4.1.}

\SGARefTarget{sga3:ii:subsection:4-3-2}%
Before continuing in this direction, let us examine more closely the
relations among \(Y\), \(\underline{\operatorname{Lie}}(Y/S)\), and
\(\underline{\operatorname{Lie}}'(Y/S)\). Let us note first that
\(\underline{\operatorname{Lie}}(\mathcal O_S/S,\mathcal M)
=\underline{\operatorname{Lie}}'(\mathcal O_S/S,\mathcal M)
=\mathbf W(\mathcal M)\) (\SGARefLink{sga3:i:subsection:4-6}{I~4.6}), so that there is a canonical
isomorphism
\[
 d:\mathcal O_S\xrightarrow{\sim}
 \underline{\operatorname{Lie}}(\mathcal O_S/S).
\]

Now let \(F\) be an \(\mathcal O_S\)-module. For every \(S'\) over \(S\)
there is a dihomomorphism
\[
 \left\{
 \begin{aligned}
   F(S)&\longrightarrow F(S'),\\
   \mathcal O(S)&\longrightarrow\mathcal O(S'),
 \end{aligned}
 \right.
\]
whence a morphism of \(\mathcal O(S')\)-modules
\[
 F(S)\otimes_{\mathcal O(S)}\mathcal O(S')\longrightarrow F(S').
\]
In particular, for \(S'=I_S(\mathcal M)\), one deduces morphisms,
functorial in \(\mathcal M\),
\[
 F(S)\otimes_{\mathcal O(S)}T_{\mathcal O_S/S}(\mathcal M)(S)
 \longrightarrow T_{F/S}(\mathcal M)(S),
\]
which, letting \(S\) vary, one may write
\[
 F\otimes_{\mathcal O_S}T_{\mathcal O_S/S}(\mathcal M)
 \longrightarrow T_{F/S}(\mathcal M).\footnotemark
\]
\footnotetext{Translator's note: the Gille--Polo re-edition (2011),
Expos\'e~II, writes the construction for two schemes \(S_2\to S_1\) over
\(S\), since both have to vary: the morphisms
\(F(S_1)\otimes_{\mathcal O(S_1)}\mathcal O(S_2)\to F(S_2)\), taken for
\(S_1=S'\) and \(S_2=I_{S'}(\mathcal M)\) with \(S'\) variable over \(S\),
form the morphism of functors written above. In the dihomomorphism the
1970 text has \(\mathcal O(S'')\) for \(\mathcal O(S')\).}
These morphisms are functorial in \(\mathcal M\), hence compatible with the
projections of the tangent bundles onto their bases. They therefore define
morphisms
\[
 F\otimes_{\mathcal O_S}
 \underline{\operatorname{Lie}}(\mathcal O_S/S,\mathcal M)
 \longrightarrow
 \underline{\operatorname{Lie}}(F/S,\mathcal M).
\]
Let us remark that these last morphisms are functorial in \(F\) and, by
construction, respect the module structures deduced from that of \(F\)
(what we have called \(\underline{\operatorname{Lie}}'(F/S,\mathcal M)\) on
the right-hand side). If \(F/S\) satisfies (E), then, since these
morphisms are also functorial in \(\mathcal M\), they also respect the
module structures deduced from that of \(\mathcal M\) by means of condition
(E), that is, those denoted \(\underline{\operatorname{Lie}}\).\footnote{Translator's
note: the Gille--Polo re-edition (2011), Expos\'e~II, deletes this
sentence, citing the example of 6.3 (\(F=\mathcal O_S\) over a field of
characteristic \(p>0\), scalars acting through \(p\)-th powers); it adds to
Definition~4.4 that the isomorphisms must respect the structures given by
(E), and states the corollary as (i)\(\Leftrightarrow\)(ii)\(\Rightarrow\)(iii),
with (ii): \(F/S\) satisfies (E) and \(d\) is an isomorphism of
\(\mathcal O_S\)-modules. The sentence is correct as printed: the operator
\(\lambda\) given by (E) acts through the homothety \(\lambda\) of
\(\mathcal M\), which induces multiplication by \(\lambda\) on
\(F\otimes\underline{\operatorname{Lie}}(\mathcal O_S/S,\mathcal M)\), so
functoriality in \(\mathcal M\) is linearity. In the example the morphisms
are zero (a scalar \(\alpha\) with \(\alpha^2=0\) acts through
\(\alpha^p=0\)): \(F\) is not good, as 6.3 says, without contradiction, and
the added requirement is automatic.}

Let us note finally that, by tensoring with
\(d:\mathcal O_S\xrightarrow{\sim}\underline{\operatorname{Lie}}(\mathcal O_S/S)\),
one deduces a morphism
\[
 F\xrightarrow{\sim}F\otimes_{\mathcal O_S}
 \underline{\operatorname{Lie}}(\mathcal O_S/S)
 \longrightarrow\underline{\operatorname{Lie}}(F/S),
\]
also denoted \(d:F\to\underline{\operatorname{Lie}}(F/S)\).\footnote{Translator's
note: the Gille--Polo re-edition (2011), Expos\'e~II, makes these morphisms
explicit (its Remarks 4.3.4 and 4.3.5). In our words: for \(f\in F(S')\)
and \(\alpha\in\mathcal O(I_{S'}(\mathcal M))\), the morphism
\(F\otimes T_{\mathcal O_S/S}(\mathcal M)\to T_{F/S}(\mathcal M)\) sends
\(f\otimes\alpha\) to \(\alpha\cdot(f\circ\rho)\), the scalar \(\alpha\)
acting on the point \(f\circ\rho\) of \(F\) over \(I_{S'}(\mathcal M)\);
restriction along the zero section multiplies by the augmentation of
\(\alpha\), so for \(\alpha\) in the augmentation ideal the image lies in
\(\underline{\operatorname{Lie}}(F/S,\mathcal M)(S')\); thus
\(d(f)=t\cdot(f\circ\rho)\) for \(\mathcal M=t\mathcal O_S\). For
\(E=\underline{\operatorname{End}}_{\mathcal O_S\text{-mod.}}(F)\), these
morphisms for \(E\) and for \(F\) are compatible with the isomorphism
\(\underline{\operatorname{Lie}}'(E/S,\mathcal M)\simeq
\underline{\operatorname{Hom}}_{\mathcal O_S\text{-mod.}}(F,\underline{\operatorname{Lie}}'(F/S,\mathcal M))\)
of 4.3, and \(E/S\) satisfies (E) if \(F/S\) does.}

\SGAInlineRefTarget{sga3:ii:definition:4-4}\subsection*{\SGARefLink{sga3:ii:definition:4-4}{Definition 4.4}.}
One says that \(F\) is a \emph{good} \(\mathcal O_S\)-module if the
morphisms
\[
 F\otimes T_{\mathcal O_S/S}(\mathcal M)\longrightarrow T_{F/S}(\mathcal M)
\]
are isomorphisms.

\SGAInlineRefTarget{sga3:ii:corollary:4-4-1}\subsection*{Corollary.}
If \(F\) is good, then
\begin{enumerate}
  \item[(i)] \(F/S\) satisfies (E);
  \item[(ii)] \(\underline{\operatorname{Lie}}(F/S,\mathcal M)
    =\underline{\operatorname{Lie}}'(F/S,\mathcal M)\);
  \item[(iii)] \(d:F\to\underline{\operatorname{Lie}}(F/S)\) is an
    isomorphism of \(\mathcal O_S\)-modules.
\end{enumerate}

Indeed: (i) It suffices to verify that
\(\underline{\operatorname{Lie}}(F/S,\mathcal M\oplus\mathcal N)
\xrightarrow{\sim}\underline{\operatorname{Lie}}(F/S,\mathcal M)\times_S
\underline{\operatorname{Lie}}(F/S,\mathcal N)\). Now if \(F\) is good,
\(F\otimes\underline{\operatorname{Lie}}(\mathcal O_S/S,\mathcal M)\to
\underline{\operatorname{Lie}}(F/S,\mathcal M)\) is an isomorphism. It only
remains to remark that the property to be proved holds for
\(\mathcal O_S\), and that
\(\underline{\operatorname{Lie}}(\mathcal O_S/S,\mathcal M)\) is isomorphic
to the product of a finite number of copies of \(\mathcal O_S\), so that
\(F\otimes\underline{\operatorname{Lie}}(\mathcal O_S/S,\mathcal M)\) is
isomorphic to the corresponding product of copies of \(F\).

(ii) The isomorphism
\(F\otimes\underline{\operatorname{Lie}}(\mathcal O_S/S,\mathcal M)\to
\underline{\operatorname{Lie}}(F/S,\mathcal M)\) respects the two module
structures defined above, and these are the same on the left-hand side.

(iii) \(d\) is the composite of two isomorphisms.\footnote{Translator's
note: the Gille--Polo re-edition (2011), Expos\'e~II, also proves a converse:
if \(F/S\) satisfies (E) and \(d\) is bijective, then \(F\) is good. In our
words: by (E) and naturality in \(\mathcal M\), the morphism
\(F\otimes\underline{\operatorname{Lie}}(\mathcal O_S/S,\mathcal M)\to
\underline{\operatorname{Lie}}(F/S,\mathcal M)\) for
\(\mathcal M=\mathcal O_S^r\) is the product of \(r\) copies of the one for
\(\mathcal M=\mathcal O_S\), which is \(d\) up to
\(F\simeq F\otimes\underline{\operatorname{Lie}}(\mathcal O_S/S)\); and on
tangent bundles the morphism is, through the zero sections and the
projections, the product of this one with the identity of \(F\).}

\SGAInlineRefTarget{sga3:ii:examples:4-4-2}\SGAInlineRefTarget{sga3:ii:example:4-4-2}\subsection*{Examples of good \(\mathcal O_S\)-modules.}
For every quasi-coherent \(\mathcal O_S\)-module \(\mathcal E\), the
\(\mathcal O_S\)-modules \(\mathbf V(\mathcal E)\) and
\(\mathbf W(\mathcal E)\) defined in \SGARefLink{sga3:i:subsection:4-6}{I~4.6} are good.\footnote{Translator's
note: the Gille--Polo re-edition (2011), Expos\'e~II, adds a proof. In our
words: for \(f\colon S'\to S\), put
\(D=\mathcal O_{S'}\oplus f^*\mathcal M\), a free \(\mathcal O_{S'}\)-module
of finite type with \(\Gamma(S',D)=\mathcal O(I_{S'}(\mathcal M))\); the
morphisms in question become
\(\operatorname{Hom}(f^*\mathcal E,\mathcal O_{S'})\otimes\Gamma(S',D)\to
\operatorname{Hom}(f^*\mathcal E,D)\) for \(\mathbf V(\mathcal E)\) and
\(\Gamma(S',f^*\mathcal E)\otimes\Gamma(S',D)\to
\Gamma(S',f^*\mathcal E\otimes D)\) for \(\mathbf W(\mathcal E)\), which are
bijective because \(D\) is a finite direct sum of copies of
\(\mathcal O_{S'}\).}

